\chapter{Linear Algebra}\label{ch:iv1:linear-algebra}

\begin{ivtrains}
Reading off rank and eigenvalues from a matrix's structure, without expansion by hand, and knowing what a covariance matrix can and cannot be. The later questions cover PCA and the SVD, conditioning, and why a sample covariance matrix misleads when there are many assets and few days.
\end{ivtrains}

\begin{question}[$\star$ \iqroles{QR, QT, DS, QD} \iqfirm{any} \iqlevel{junior}]\label{iq:iv1:linear-algebra:1}
Let $u = (1, 2, 3)$ and $v = (1, 1, 1)$. What are the rank and the eigenvalues of the matrix $uv^\top$? What are the eigenvalues of the $n$ by $n$ matrix of all ones, and of the matrix with zeros on the diagonal and ones elsewhere?
\end{question}

\begin{question}[$\star$ \iqroles{QR, QT, DS, Risk} \iqfirm{any} \iqlevel{junior}]\label{iq:iv1:linear-algebra:2}
Ten assets have unit variance and a common pairwise correlation of $\rho = 0.2$. What are the eigenvalues of the correlation matrix, and what are the eigenvectors? What is the lowest possible value of $\rho$ for 10 assets, and what is the variance of their equally weighted average for large numbers?
\end{question}

\begin{question}[$\star$ \iqroles{QR, QT, Risk, DS} \iqfirm{any} \iqlevel{junior}]\label{iq:iv1:linear-algebra:3}
Three assets have pairwise correlations $\rho_{12} = 0.9$, $\rho_{23} = 0.9$ and $\rho_{13} = 0$, as given by a data vendor. Is this a valid correlation matrix? Prove your answer by finding a portfolio with an impossible variance.
\end{question}

\begin{question}[$\star$ \iqroles{QR, QT, DS, QD} \iqfirm{any} \iqlevel{junior}]\label{iq:iv1:linear-algebra:4}
A matrix $P$ has the property $P^2 = P$. What can its eigenvalues be? Show that its rank equals its trace. What is the trace of the matrix $X(X^\top X)^{-1}X^\top$, where $X$ has $n$ rows and $p$ independent columns?
\end{question}

\begin{question}[$\star$ \iqroles{QR, QT, QD, DS} \iqfirm{any} \iqlevel{junior}]\label{iq:iv1:linear-algebra:5}
Let $A$ be the 2 by 2 symmetric matrix with 2 on the diagonal and 1 off it. Find its eigenvalues and eigenvectors, and use them to write down $A^{10}$ and $A^{-1}$ exactly, with no repeated multiplication.
\end{question}

\begin{question}[$\star\star$ \iqroles{QR, QT, Risk, DS} \iqfirm{any} \iqlevel{junior}]\label{iq:iv1:linear-algebra:6}
Asset $X$ has a correlation of $0.9$ with $Y$, and $Y$ has a correlation of $0.9$ with $Z$. What are the lowest and highest possible correlations between $X$ and $Z$? Derive them in two ways: from a determinant and from angles between vectors.
\end{question}

\begin{question}[$\star\star$ \iqroles{QR, QT, QD} \iqfirm{any} \iqlevel{junior}]\label{iq:iv1:linear-algebra:7}
Let $P$ be the $n$ by $n$ cyclic shift, with $(Px)_i = x_{i + 1}$, indices wrapping around. What are its eigenvalues and eigenvectors, and what is its condition number? Use them to find the eigenvalues of the 4 by 4 circulant whose first row is $(2, 1, 0, 0)$.
\end{question}

\begin{question}[$\star\star$ \iqroles{QR, QD, QT} \iqfirm{any} \iqlevel{junior}]\label{iq:iv1:linear-algebra:8}
Take the $n$ by $n$ tridiagonal matrix with 2 on the diagonal and $-1$ just off it. Find its eigenvalues for $n = 3$, then give the general formula. What happens to the condition number as $n$ grows, and why does it matter when solving a finite-difference problem?
\end{question}

\begin{question}[$\star\star$ \iqroles{QR, QT, DS, Risk} \iqfirm{hedge fund} \iqlevel{junior}]\label{iq:iv1:linear-algebra:9}
Two assets have a correlation of $0.5$. If both have unit volatility, what are the two principal components and their shares of variance? If the first asset's volatility is doubled, how much variance does the first component explain? Why does this show that PCA depends on units, and what do you do about it?
\end{question}

\begin{question}[$\star\star$ \iqroles{QR, QD, DS} \iqfirm{any} \iqlevel{junior}]\label{iq:iv1:linear-algebra:10}
Find the singular values of the matrix with rows $(3, 0)$ and $(4, 5)$, and its condition number in the 2-norm. Compare them with the eigenvalues of the matrix. How do singular values relate to the principal components of a data matrix?
\end{question}

\begin{question}[$\star\star$ \iqroles{QR, QD, QT, DS} \iqfirm{any} \iqlevel{junior}]\label{iq:iv1:linear-algebra:11}
Solve $Ax = b$ for $A$ with rows $(1, 1)$ and $(1, 1.0001)$, and for $b = (2, 2)$ and then $b = (2, 2.0001)$. What does the result show? Give the condition number of $A$, and the rule of thumb linking it to lost digits. What does squaring $A$, as in the normal equations of a regression, do?
\end{question}

\begin{question}[$\star\star$ \iqroles{QR, QT, DS, Risk} \iqfirm{hedge fund, bank} \iqlevel{junior}]\label{iq:iv1:linear-algebra:12}
You compute the sample covariance matrix of 500 stocks from 250 daily returns. What is the largest its rank can be? What follows for a minimum-variance portfolio computed from it? Give two practical repairs.
\end{question}

\begin{question}[$\star\star$ \iqroles{QR, QT, Risk, DS} \iqfirm{hedge fund} \iqlevel{junior}]\label{iq:iv1:linear-algebra:13}
A covariance matrix has eigenvalues 5, 1 and $0.1$. You replace it with $0.7$ times itself plus $0.3$ times $\mu I$, with $\mu$ the mean of the eigenvalues. What are the new eigenvalues and the new condition number? What stays the same? Why does this help an optimiser, and where would $0.3$ come from?
\end{question}

\begin{question}[$\star\star$ \iqroles{QR, QT, Risk} \iqfirm{hedge fund} \iqlevel{junior}]\label{iq:iv1:linear-algebra:14}
One asset has a volatility of 10\%, another of 20\%, and their correlation is $\rho$. Find the minimum-variance weights. For what $\rho$ is the second asset's weight exactly zero? What are the weights and the portfolio's volatility at $\rho = 0.8$? Why is the variance below that of the first asset, when the second is riskier and correlated?
\end{question}

\begin{question}[$\star\star$ \iqroles{QR, QT, DS, QD} \iqfirm{any} \iqlevel{junior}]\label{iq:iv1:linear-algebra:15}
A market state moves between two regimes: from calm, it stays calm with chance $0.9$; from stressed, it stays stressed with chance $0.8$. Write the transition matrix, its eigenvalues and its stationary distribution. How quickly does the effect of today's state die away? What is that time for a half-way decay?
\end{question}

\begin{question}[$\star\star$ \iqroles{QR, QD, Dev} \iqfirm{any} \iqlevel{junior}]\label{iq:iv1:linear-algebra:16}
Write the Fibonacci recursion as a 2 by 2 matrix. What are its eigenvalues? Derive the closed form, and say how many multiplications matrix powers need for $F_n$. Using double-precision floating point, does the closed form always give the exact integer?
\end{question}

\begin{question}[$\star\star$ \iqroles{QR, QT, QD} \iqfirm{any} \iqlevel{junior}]\label{iq:iv1:linear-algebra:17}
A 100 by 100 matrix has entries $a_{ij} = i + j$, with indices from 1 to 100. What is its rank? Find its nonzero eigenvalues in closed form, without expanding anything, by using the structure of the matrix.
\end{question}

\begin{question}[$\star\star\star$ \iqroles{QR, QT, DS, Risk} \iqfirm{hedge fund} \iqlevel{junior, experienced}]\label{iq:iv1:linear-algebra:18}
Five hundred uncorrelated assets with unit variance are observed for 1000 days. Between which values do the sample covariance eigenvalues lie, by the Marchenko--Pastur law, although all the true eigenvalues are 1? A minimum-variance portfolio is built from the sample matrix. How does its risk in the sample compare with its risk out of sample, and with the true optimum?
\end{question}

\begin{question}[$\star\star\star$ \iqroles{QR, QT, QD, Risk} \iqfirm{hedge fund} \iqlevel{junior, experienced}]\label{iq:iv1:linear-algebra:19}
A factor risk model has covariance matrix $\Sigma = D + \beta\beta^\top$ for 1000 assets, with $D$ diagonal. Find $\Sigma^{-1}$ with the Sherman--Morrison formula, and prove it. What is the cost compared with a direct inverse? Write the minimum-variance weights, and say what they do to high-beta stocks.
\end{question}

\begin{question}[$\star\star\star$ \iqroles{QR, QT, DS} \iqfirm{hedge fund} \iqlevel{junior, experienced}]\label{iq:iv1:linear-algebra:20}
$X$, $Y$ and $Z$ have unit variance and all pairwise correlations equal to $0.5$. What is the partial correlation of $X$ and $Y$ given $Z$? Compute it in two ways: from residuals of regressions on $Z$, and from the inverse of the correlation matrix. What is the partial correlation in general if the correlation is $\rho$?
\end{question}

\begin{question}[$\star\star\star$ \iqroles{QR, QD, DS} \iqfirm{any} \iqlevel{junior, experienced}]\label{iq:iv1:linear-algebra:21}
You need the top eigenvector of a 1000 by 1000 covariance matrix, whose two top eigenvalues are 10 and 5. How does power iteration work, and about how many steps does it take for an error of $10^{-6}$ in the vector? And in the eigenvalue? How many steps if the second eigenvalue is $9.9$? When could it fail?
\end{question}

\begin{question}[$\star\star\star$ \iqroles{QR, QD, DS, QT} \iqfirm{hedge fund} \iqlevel{junior, experienced}]\label{iq:iv1:linear-algebra:22}
Show that $AB$ and $BA$ have the same nonzero eigenvalues, even if they have different sizes. Use this to compute the principal components of a $T$ by $N$ return matrix with $T = 250$ and $N = 5000$, without forming the $N$ by $N$ matrix. What is the eigenvector map? Also derive $\det(I + uv^\top) = 1 + v^\top u$.
\end{question}

\begin{question}[$\star\star\star$ \iqroles{QR, QT, Risk, QD} \iqfirm{bank, hedge fund} \iqlevel{junior, experienced}]\label{iq:iv1:linear-algebra:23}
A correlation matrix built from a mix of sources has $\rho_{12} = \rho_{23} = 0.9$ and $\rho_{13} = 0$. Give its eigenvalues. Repair it by clipping the negative eigenvalue to zero and rescaling to a unit diagonal, and give the new correlations. In what sense is this close to the original, and is it the closest possible?
\end{question}

\begin{question}[$\star\star\star\star$ \iqroles{QR, QT} \iqfirm{hedge fund} \iqlevel{experienced}]\label{iq:iv1:linear-algebra:24}
An $n$ by $n$ matrix has independent entries, each with mean 0 and variance 1. Find $E[\det(A)^2]$ exactly. Check the answer for $n = 3$ with $\pm 1$ entries, by reasoning about the size of the enumeration. What does this tell you about how the typical size of a determinant grows? Compare it with the largest possible determinant of a $\pm 1$ matrix.
\end{question}

\begin{question}[$\star\star\star\star$ \iqroles{QR, QT, QD} \iqfirm{hedge fund} \iqlevel{experienced}]\label{iq:iv1:linear-algebra:25}
Prove that the element-by-element product of two positive semi-definite matrices is itself positive semi-definite. Use this to show that multiplying a correlation matrix, entry by entry, by a matrix with entries $e^{-\lambda |i - j|}$ gives a valid correlation matrix. Also prove the bound $\lambda_{\min}(A \circ B) \ge \lambda_{\min}(A) \min_i B_{ii}$.
\end{question}
